\vspace{2.5pt plus 3pt minus 0.5pt}\noindent\hspace*{1.2em}\kbvnum{(4)}\hspace{0.55em}ಕೆಳಗೆ ಕೊಟ್ಟಿರುವ ಪದಗಳಲ್ಲಿ ಕಾಣಸಿಗುವ ವರ್ಗೀಯ ವ್ಯಂಜನ, ಅವರ್ಗೀಯ ವ್ಯಂಜನ ಮತ್ತು ಯೋಗವಾಹಗಳನ್ನು ಗುರ್ತಿಸಿ:

\vspace{5pt plus 4pt minus 1pt}
\noindent{\renewcommand{\arraystretch}{1.24}\setlength{\tabcolsep}{4pt}
\begin{tabularx}{\textwidth}{X X X X X}
ಅಗಸ & ಅಂಗಳ & ಆಸನ & ಆಭರಣ & ಔನಃ \\
ಉರಗ & ಓಲಗ & ಐದಳೆ & ಮದಗಜ & ಆಂತಃ \\
\end{tabularx}}
\vspace{5pt plus 4pt minus 1pt}

\vspace{2.5pt plus 3pt minus 0.5pt}\noindent\hspace*{1.2em}\kbvnum{(5)}\hspace{0.55em}ಗುಣಿತಾಕ್ಷರ ಎಂದರೇನು? “ಚ್” ವ್ಯಂಜನದ ಗುಣಿತಾಕ್ಷರಗಳ ಬರೆದು ತಿಳಿಸಿ.\par

\vspace{2.5pt plus 3pt minus 0.5pt}\noindent\hspace*{1.2em}\kbvnum{(6)}\hspace{0.55em}ಸಜಾತೀಯ ಒತ್ತಕ್ಷರವುಳ್ಳ ನಾಲ್ಕು ಪದಗಳನ್ನು ಬರೆಯಿರಿ.\par

\vspace{2.5pt plus 3pt minus 0.5pt}\noindent\hspace*{1.2em}\kbvnum{(7)}\hspace{0.55em}ವಿಜಾತೀಯ ಒತ್ತಕ್ಷರವುಳ್ಳ ನಾಲ್ಕು ಪದಗಳನ್ನು ಬರೆಯಿರಿ.\par

\vspace{2.5pt plus 3pt minus 0.5pt}\noindent\hspace*{1.2em}\kbvnum{(8)}\hspace{0.55em}ಗ್ಗ, ಜ್ರ, ನ್ಯ, ಶ್ಚ, ರೃ, ಸ್ಹ, ದೃ, ಜ್ವ—ಈ ಒತ್ತ ಗನ್ನು ಬಿಡಿಸಿ ಬರೆಯಿರಿ.\par

\vspace{2.5pt plus 3pt minus 0.5pt}\noindent\hspace*{1.2em}\kbvnum{(9)}\hspace{0.55em}“ಸೀತಾಪತಿ” -ಈ ಶಬ್ದದಲ್ಲಿರುವ ಸ್ವರ, ವ್ಯಂಜನಗಳ ಬಿಡಿಸಿ ಬರೆಯಿರಿ.\par

\vspace{2.5pt plus 3pt minus 0.5pt}\noindent\hspace*{1.2em}\kbvnum{(10)}\hspace{0.55em}ಕ ಕಾರದಿಂದ ಳ ಕಾರದವರೆಗಿನ ಎಲ್ಲ ವ್ಯಂಜನಗಳ ಗುಣಿತಾಕ್ಷರಗಳನ್ನೊಳಗೊಂಡ ಪಟ್ಟಿ ಸಿದ್ಧಪಡಿಸಿ (ಯಾವುದೇ ಅಕ್ಷರವನ್ನು ಹೆಸರಿಸುವಾಗ ಅದರ ಮುಂದೆ “ಕಾರ” ಎಂಬ ಪ್ರತ್ಯಯ ಸೇರಿಸುವುದು ರೂಢಿಯಲ್ಲಿದೆ.)\par

\kbvchapter{ಪಾಠ ೩}{ನಾಮಪದ}

\noindent\hspace*{2.2em}ಕೋತಿ ಲಾಗ ಹಾಕಿತು.\par

\noindent\hspace*{2.2em}ಮೇಜು ಬಹಳ ದೊಡ್ಡದು.\par

\noindent\hspace*{2.2em}ಬೆಂಗಳೂರು ದೊಡ್ಡ ವಾಣಿಜ್ಯ ಕೇಂದ್ರ.\par

\noindent\hspace*{2.2em}ರಾಮಯ್ಯ ಈ ಕಡೆ ಬರುತ್ತಿದ್ದಾನೆ.\par

\indent ಈ ವಾಕ್ಯಗಳಲ್ಲಿ ದಪ್ಪ ಅಕ್ಷರಗಳಲ್ಲಿರುವ ಪದಗಳನ್ನು ಚೆನ್ನಾಗಿ ಗಮನಿಸಿ-- ಕೋತಿ - ಒಂದು ಪ್ರಾಣಿ ; ಮೇಜು - ಒಂದು ವಸ್ತು ; ಬೆಂಗಳೂರು - ಒಂದು  ನಗರ ; ರಾಮಯ್ಯ - ಒಬ್ಬ ವ್ಯಕ್ತಿ. ಹೀಗೆ ಪ್ರಾಣಿ, ವಸ್ತು, ಸ್ಥಳ, ಮನುಷ್ಯ ಇವೇ ಮೊದಲಾದವುಗಳನ್ನು ಹೆಸರಿಸುವ ಪದಗಳಿಗೆ “ನಾಮ ಪ್ರಕೃತಿ”ಗಳೆಂದು  ಹೆಸರು.\par
